% Actual class: 02
% Slide decks: Course Introduction - S2; Part 1 (Intro and Background); Tutorial One
% Exact scope: Course Introduction pp. 1--10 (administrative); Part 1 pp. 1--46; Tutorial One p. 1
% Video supplement: Part 1.1, Part 1.2, Part 1.3 (complete)
% Human verified: Yes

\section{第 02 节课：OS Foundations、Interrupts 与 System Calls}
\label{lecture:SC2005-L02}

\lecturemeta
  {SC2005-L02}
  {2026-08-14}
  {Course Introduction - S2；Part 1 (Intro and Background)；Tutorial One}
  {Course Introduction pp.~1--10（课程行政信息）；Part 1 pp.~1--46；Tutorial One p.~1}
  {Part 1.1、Part 1.2、Part 1.3 全部视频}
  {已确认（2026-08-14；讲解与问答核验）}

\subsection{本讲主线}

Operating System（操作系统，OS）用 abstraction（抽象）隐藏硬件细节，并在多个 programs（程序）之间分配 CPU、memory（内存）和 I/O devices（输入输出设备）。Multiprogramming（多道程序设计）依靠 interrupt（中断）在 CPU 工作与 I/O 等待之间重叠执行；dual-mode operation（双模式运行）和 memory protection（内存保护）限制普通程序的权限；system call（系统调用）则提供从 user mode（用户态）进入 kernel mode（内核态）的受控入口。

\lecturetopic{SC2005-L02-T01}{Operating System 的角色与目标}

\keyidea{OS 是 application（应用）与 hardware（硬件）之间的管理层：它提供抽象、分配资源并执行保护策略，但普通用户指令通常直接在 CPU 上运行，并非每条指令都要经过 OS。}

\begin{center}
\small
\begin{tabularx}{\textwidth}{@{}>{\raggedright\arraybackslash}p{3.1cm}>{\raggedright\arraybackslash}p{4.4cm}X@{}}
  \toprule
  视角 & 核心含义 & 代表性抽象或职责 \\
  \midrule
  Abstraction provider（抽象提供者） & 隐藏 hardware complexity（硬件复杂性），提供稳定 interface（接口） & file（文件）、process（进程）、virtual memory（虚拟内存） \\
  Resource allocator（资源分配器） & 在竞争者之间分配有限资源 & CPU scheduling（CPU 调度）、memory allocation（内存分配）、I/O scheduling（I/O 调度） \\
  Control program（控制程序） & 检查操作是否合法，防止程序相互破坏或破坏 OS & privilege（特权）、access control（访问控制）、fault handling（故障处理） \\
  Kernel（内核） & OS 中以高 privilege level（特权级）运行的核心部分 & 响应 system call、interrupt 与 exception（异常） \\
  \bottomrule
\end{tabularx}
\end{center}

OS 同时追求 user convenience（用户便利性）和 efficient hardware utilization（高效利用硬件）。精细的资源分配通常需要更多 workload information（工作负载信息），而要求用户提供过多信息又会降低便利性，因此两项目标之间存在 trade-off（权衡）。

\lecturetopic{SC2005-L02-T02}{Multiprogramming、Time-sharing 与 Multiprocessing}

\begin{center}
\small
\begin{tabularx}{\textwidth}{@{}>{\raggedright\arraybackslash}p{3.0cm}>{\raggedright\arraybackslash}p{4.4cm}X@{}}
  \toprule
  系统/机制 & 核心机制 & 首要目的 \\
  \midrule
  Simple batch system（简单批处理系统） & memory 中通常只有一个 user job；该 job 等待 I/O 时 CPU 可能空闲 & 简化自动作业序列 \\
  Multiprogramming（多道程序设计） & 多个 jobs 同驻 main memory；一个 job 阻塞时运行另一个 & 提高 CPU utilization（利用率） \\
  Time-sharing（分时） & 在 multiprogramming 基础上快速切换并允许 preemption（抢占） & 提供 interactive response（交互响应） \\
  Multiprocessing / multicore（多处理/多核） & 多个 cores（核心）可同时执行 instruction streams（指令流） & 获得 true parallelism（真正并行）与更高 throughput（吞吐量） \\
  \bottomrule
\end{tabularx}
\end{center}

\begin{figure}[htbp]
  \centering
  \resizebox{0.96\textwidth}{!}{\begin{tikzpicture}[
  >=Latex,
  title/.style={font=\bfseries,align=center},
  lane/.style={draw=noteBlue,rounded corners=2pt,minimum height=0.72cm,align=center},
  flow/.style={->,thick,draw=noteBlue}
]
  \node[title] at (3.45,2.0) {Single core\\concurrency（并发）};
  \node[anchor=east] at (0.75,0.85) {CPU};
  \node[lane,fill=blue!8,minimum width=1.9cm] at (1.75,0.85) {Job A};
  \node[lane,fill=red!8,minimum width=1.9cm] at (3.65,0.85) {Job B};
  \node[lane,fill=blue!8,minimum width=1.9cm] at (5.55,0.85) {Job A};
  \draw[flow] (0.8,-0.1) -- (6.55,-0.1);
  \node[anchor=west,font=\small] at (6.6,-0.1) {time};

  \draw[dashed,gray] (7.4,-0.55) -- (7.4,2.35);

  \node[title] at (11.25,2.0) {Two cores\\parallelism（并行）};
  \node[anchor=east] at (8.65,1.05) {Core 0};
  \node[lane,fill=blue!8,minimum width=4.9cm] at (11.25,1.05) {Job A};
  \node[anchor=east] at (8.65,0.2) {Core 1};
  \node[lane,fill=red!8,minimum width=4.9cm] at (11.25,0.2) {Job B};
  \draw[flow] (8.75,-0.65) -- (13.75,-0.65);
  \node[anchor=west,font=\small] at (13.8,-0.65) {time};
\end{tikzpicture}
}
  \caption{单核 multiprogramming（多道程序设计）通过切换形成 concurrency（并发）；multicore hardware（多核硬件）允许多个执行流真正 parallel（并行）。}
  \label{fig:sc2005-l02-concurrency-parallelism}
\end{figure}

Multiprogramming 是 OS mechanism（操作系统机制），multiprocessing 是 hardware capability（硬件能力）。Multithreading（多线程）描述一个 process 内存在多个 threads（线程）；单核上也能并发执行多个 threads，因此它不能作为 multiprocessing 的同义词。

Real-time system（实时系统）的正确性同时取决于计算结果和完成时间：关键任务必须在 deadline（截止时间）前完成。它不等于“平均运行得很快”；general-purpose OS（通用操作系统）偶尔错过 deadline，仍不能据此称为 hard real-time OS（硬实时操作系统）。

\textbf{对应练习：}\relatedexercise{SC2005-TUT01-Q03}{Multiprogramming 与 Multiprocessing}。

\lecturetopic{SC2005-L02-T03}{Interrupt-driven I/O、Context 与 DMA}

CPU、main memory 和 device controllers（设备控制器）通过 bus（总线）连接。每个 controller 通常管理特定设备并带有 local buffer（本地缓冲区）；CPU 可以执行程序，同时 controller 独立处理较慢的 I/O request（I/O 请求）。

\begin{figure}[htbp]
  \centering
  \resizebox{\textwidth}{!}{\begin{tikzpicture}[
  >=Latex,
  cpu/.style={draw=noteBlue,rounded corners=2pt,fill=blue!7,minimum height=0.8cm,align=center},
  dev/.style={draw=ntuRed,rounded corners=2pt,fill=red!6,minimum height=0.8cm,align=center},
  flow/.style={->,thick,draw=noteBlue},
  signal/.style={->,very thick,draw=ntuRed},
  stepblue/.style={circle,fill=noteBlue,text=white,inner sep=1.5pt,font=\scriptsize\bfseries},
  stepred/.style={circle,fill=ntuRed,text=white,inner sep=1.5pt,font=\scriptsize\bfseries}
]
  \node[font=\bfseries,anchor=east] at (1.35,1.2) {CPU lane};
  \node[cpu,minimum width=1.9cm] (pa) at (2.5,1.2) {Process A};
  \node[cpu,minimum width=2.1cm] (os1) at (5.0,1.2) {OS: start I/O};
  \node[cpu,minimum width=2.3cm] (pb) at (7.9,1.2) {Process B};
  \node[cpu,minimum width=1.5cm] (isr) at (10.45,1.2) {ISR};
  \node[cpu,minimum width=2.3cm] (next) at (12.9,1.2) {next process};

  \node[font=\bfseries,anchor=east] at (1.35,-0.9) {Device lane};
  \node[dev,minimum width=4.8cm] (transfer) at (7.9,-0.9) {controller transfers data};

  \draw[flow] (pa) -- node[stepblue]{1} (os1);
  \draw[flow] (os1) -- (pb);
  \draw[flow] (pb) -- (isr);
  \draw[flow] (isr) -- node[stepblue]{4} (next);
  \draw[signal] (os1.south) -- node[stepred]{2} (transfer.west);
  \draw[signal] (transfer.east) -- node[stepred]{3} (isr.south);
\end{tikzpicture}
}
  \caption{Interrupt-driven I/O（中断驱动 I/O）：1，Process A 通过 system call 进入 OS；2，OS 启动 I/O；3，device 完成后产生 hardware interrupt；4，ISR 结束后由 OS 继续调度。}
  \label{fig:sc2005-l02-interrupt-timeline}
\end{figure}

典型过程为：Process A 发起 system call；OS 配置 controller 并阻塞 A；scheduler（调度器）运行 Process B；设备完成传输后产生 hardware interrupt（硬件中断）；CPU 保存必要状态并执行 ISR（Interrupt Service Routine，中断服务程序）；OS 再决定恢复 A、B 或其他 process。

\begin{center}
\small
\begin{tabularx}{0.94\textwidth}{@{}>{\raggedright\arraybackslash}p{3.3cm}>{\raggedright\arraybackslash}p{3.1cm}X@{}}
  \toprule
  事件 & 来源与时序 & 典型用途 \\
  \midrule
  Hardware interrupt（硬件中断） & 外部 device 产生，asynchronous（异步） & I/O 完成、timer（定时器）到期 \\
  Trap / synchronous exception（陷阱/同步异常） & 当前 CPU instruction 引起，synchronous（同步） & system call、除零、非法地址、非法特权指令 \\
  \bottomrule
\end{tabularx}
\end{center}

Interrupt entry（中断入口）必然需要 context save（上下文保存），但不必然产生 process context switch（进程上下文切换）。只有 scheduler 决定改运行另一个 process 时才发生后者；若 ISR 结束后返回同一 process，则只经历 privilege/control transfer（特权级/控制流切换）。处理 ISR 时也不是永久禁止全部中断：实现通常短暂 mask（屏蔽）部分中断，并可允许更高优先级中断嵌套。

\begin{figure}[htbp]
  \centering
  \resizebox{0.94\textwidth}{!}{\begin{tikzpicture}[
  >=Latex,
  box/.style={draw=noteBlue,rounded corners=2pt,fill=blue!6,minimum width=3.0cm,minimum height=1.0cm,align=center},
  flow/.style={->,thick,draw=noteBlue},
  setup/.style={->,dashed,thick,draw=ntuRed}
]
  \node[font=\bfseries,anchor=east] at (2.9,1.5) {Ordinary I/O（课程抽象）};
  \node[box] (d1) at (5.0,1.5) {device controller\\local buffer};
  \node[box] (c1) at (9.2,1.5) {CPU / driver\\参与复制};
  \node[box] (m1) at (13.4,1.5) {main memory};
  \draw[flow] (d1) -- (c1);
  \draw[flow] (c1) -- (m1);

  \node[font=\bfseries,anchor=east] at (2.9,-0.8) {DMA（直接内存访问）};
  \node[box] (cpu) at (5.0,-0.8) {CPU / OS};
  \node[box] (d2) at (9.2,-0.8) {DMA-capable\\controller};
  \node[box] (m2) at (13.4,-0.8) {main memory};
  \draw[setup] (cpu) -- node[above,font=\small]{setup} (d2);
  \draw[flow] (d2) -- (m2);
  \draw[setup] ([yshift=-0.18cm]d2.south west) to[bend left=24]
    node[below,font=\small]{completion interrupt} ([yshift=-0.18cm]cpu.south east);
\end{tikzpicture}
}
  \caption{课程抽象中的 ordinary interrupt-driven I/O（普通中断驱动 I/O）与 DMA（直接内存访问）。}
  \label{fig:sc2005-l02-dma-paths}
\end{figure}

对于 ordinary interrupt-driven I/O，完成中断到达时，数据可能仍在 controller buffer，随后由 CPU/driver（驱动）参与搬入 main memory。对于 DMA（Direct Memory Access，直接内存访问），OS 先设置 memory region（内存区域）和 transfer count（传输计数），controller 再直接传输整个 block（数据块），完成后通常只产生一次 interrupt。

“Without CPU intervention（无需 CPU 介入）”仅指 bulk data movement（批量数据搬运）阶段；CPU 仍负责配置 DMA 和处理完成中断。“每 byte 一个中断、每 block 一个中断”是课程用于对比的简化模型，实际设备也可能进行 buffering（缓冲）和 batching（批处理）。

Storage hierarchy（存储层次）按典型距离可概括为 registers $\rightarrow$ cache $\rightarrow$ main memory $\rightarrow$ secondary storage（寄存器、缓存、主存、二级存储）。越接近 CPU 通常越快、越小、单位成本越高；本讲仅作背景，不作为当前考核重点。

\textbf{对应练习：}\relatedexercise{SC2005-TUT01-Q01}{Interrupt-driven I/O 与 DMA}。

\lecturetopic{SC2005-L02-T04}{Dual-mode、I/O 与 Memory Protection}

Dual-mode operation 用 hardware privilege（硬件特权级）区分 user mode，以及 kernel mode（内核态）、monitor mode（监控态）或 supervisor mode（管理态）。普通程序在 user mode 执行；interrupt 或 trap 使 CPU 进入 kernel mode；return-from-trap（陷阱返回）再恢复 user mode。

Root/admin（管理员账户）是 OS 内部的 account/permission concept（账户/权限概念），不是 CPU mode。管理员启动的普通程序仍通常在 user mode，只能通过 system call 或获授权的 kernel extension（内核扩展）间接影响 kernel-mode code。

课程模型将 I/O instructions（I/O 指令）视为 privileged instructions（特权指令）：普通程序必须请求 OS 检查设备、操作和访问权限。现代系统也可能使用 memory-mapped I/O（内存映射 I/O），但相应 address range（地址范围）仍由 privilege 与 memory mapping（内存映射）保护。

在 base-limit protection（基址-界限保护）的简化模型中，base register 保存第一个合法 physical address（物理地址），limit register 保存合法区域长度。访问合法当且仅当
\[
  \boxed{\text{base}\leq \text{address}<\text{base}+\text{limit}}.
\]
所以最后一个合法 byte address 为 $\text{base}+\text{limit}-1$。加载 base/limit 的指令本身必须是 privileged；否则程序可以自行扩大可访问区域。

\textbf{对应练习：}\relatedexercise{SC2005-TUT01-Q02}{Protection 与 Trap 判断}。

\lecturetopic{SC2005-L02-T05}{OS Services、API 与 System Call}

OS services（操作系统服务）包括 program execution（程序执行）、I/O、file systems（文件系统）、communication（通信）、resource allocation（资源分配）、accounting（统计）、error detection（错误检测）以及 protection/security（保护与安全）。System call 是 user program 请求这些 kernel services（内核服务）的受控接口。

\begin{figure}[htbp]
  \centering
  \resizebox{0.88\textwidth}{!}{\begin{tikzpicture}[
  >=Latex,
  user/.style={draw=noteBlue,rounded corners=2pt,fill=blue!6,minimum width=3.5cm,minimum height=0.95cm,align=center},
  kernel/.style={draw=ntuRed,rounded corners=2pt,fill=red!6,minimum width=3.5cm,minimum height=0.95cm,align=center},
  hardware/.style={draw=black!65,rounded corners=2pt,fill=black!5,minimum width=3.5cm,minimum height=0.95cm,align=center},
  down/.style={->,very thick,draw=noteBlue},
  up/.style={->,very thick,draw=ntuRed},
  stepblue/.style={circle,fill=noteBlue,text=white,inner sep=1.5pt,font=\scriptsize\bfseries},
  stepred/.style={circle,fill=ntuRed,text=white,inner sep=1.5pt,font=\scriptsize\bfseries}
]
  \node[font=\bfseries,anchor=east] at (1.4,2.15) {User mode};
  \node[user] (app) at (4.6,2.15) {application / library API};

  \draw[thick] (0.8,1.15) -- (13.2,1.15);
  \node[anchor=east,font=\small,fill=white,inner sep=1pt] at (13.05,1.38)
    {user/kernel privilege boundary};

  \node[font=\bfseries,anchor=east] at (1.4,0.0) {Kernel mode};
  \node[kernel] (entry) at (4.6,0.0) {system-call handler};
  \node[kernel] (isr) at (10.4,0.0) {device driver / ISR};

  \draw[thick] (0.8,-1.05) -- (13.2,-1.05);
  \node[font=\bfseries,anchor=east] at (1.4,-2.05) {Hardware};
  \node[hardware] (dev) at (10.4,-2.05) {I/O device controller};

  \draw[down] (app) -- node[stepblue]{1} (entry);
  \draw[up] (dev) -- node[stepred]{2} (isr);
\end{tikzpicture}
}
  \caption{两种进入 kernel 的方向：1，system call/trap（系统调用/陷阱）由当前程序同步触发；2，hardware interrupt（硬件中断）由 device 异步触发。}
  \label{fig:sc2005-l02-privilege-boundary}
\end{figure}

Library API（库接口）不等于 system call。以 C 为例，\texttt{fopen()} 是 standard library function（标准库函数）；其实现可能调用 \texttt{open()} system call。Library wrapper（库封装）准备 system-call number（系统调用号）和 arguments（参数），再执行专门的 \texttt{syscall}/\texttt{svc} instruction；hardware 据此切换 privilege level 并跳转到注册的 kernel entry point（内核入口），而不是让 user program 普通跳转到 kernel memory。

并非所有 API calls 都进入 kernel：例如只在 user space（用户空间）计算字符串长度的 \texttt{strlen()} 通常不产生 system call。反过来，一个 API call 也可能触发多个 system calls。即使调用了 file I/O system call，若数据已在 page cache（页缓存）中，也不一定产生新的 physical-device interrupt（物理设备中断）。

\subsection{一分钟回顾}
\par\noindent
OS 通过 abstractions（抽象）提供易用接口，通过 scheduling（调度）和 multiprogramming 提高资源利用率，并用 privilege 与 address checks（地址检查）保护系统。System call/trap 是 program 到 kernel 的同步受控入口；hardware interrupt 是 device 到 CPU 的异步通知。DMA 让 controller 直接搬运 data block，但 CPU 仍负责配置和完成处理。Concurrency（并发）可以存在于单核系统，parallelism（并行）则要求多个执行资源。
